\documentclass[12pt,a4paper]{article}

\usepackage[a4paper,margin=2.1cm]{geometry}
\usepackage[T2A]{fontenc}
\usepackage[utf8]{inputenc}
\usepackage[russian]{babel}
\usepackage{amsmath,amssymb,mathtools}
\usepackage{array,booktabs,longtable}
\usepackage{xcolor}
\usepackage{hyperref}
\usepackage{tikz}

\hypersetup{
  colorlinks=true,
  linkcolor=blue!60!black,
  urlcolor=blue!60!black
}

\setlength{\parindent}{1.25cm}
\setlength{\parskip}{0.25em}
\renewcommand{\arraystretch}{1.25}

\begin{document}

\begin{titlepage}
\thispagestyle{empty}
\centering

\vspace*{0.22\textheight}

{\Huge\bfseries Ревью домашней работы\par}

\vspace{2.2cm}

{\Large Автор: Лёвин Артём Александрович\par}

\vspace{0.8cm}

{\large 20 сентября 2026 года\par}

\vfill

\begin{tikzpicture}
\draw[blue!35,line width=0.9pt]
  (0,0) .. controls (3,0.8) and (7,-0.8) .. (11,0);
\end{tikzpicture}

\end{titlepage}

\newpage
\section*{Сводная таблица}
\small
\setlength{\tabcolsep}{3pt}
\begin{longtable}{>{\centering\arraybackslash}p{0.8cm} p{2.1cm} p{4.5cm} p{2.5cm} p{3.8cm}}
\toprule
\textbf{№} & \textbf{Тема} & \textbf{Суть ошибки} & \textbf{Ваш ответ} & \textbf{Правильный ответ} \\
\midrule
\endfirsthead
\toprule
\textbf{№} & \textbf{Тема} & \textbf{Суть ошибки} & \textbf{Ваш ответ} & \textbf{Правильный ответ} \\
\midrule
\endhead
\bottomrule
\endfoot

\hyperref[task:3]{3} &
Производная и касательная &
Решение не выполнено: не использован геометрический смысл производной как углового коэффициента касательной. &
не указан &
$-\dfrac15$ \\

\hyperref[task:7]{7} &
Кредит с ежегодным начислением процентов &
Решение не выполнено; не составлена рекуррентная схема изменения долга при трёх равных платежах. &
не указан &
$259\,200$ руб. \\

\hyperref[task:10]{10} &
Максимум мощности нагрузки &
Получена верная формула мощности, но не выполнен шаг, позволяющий найти её максимальное значение. &
не указан &
$625$ Вт \\

\hyperref[task:11]{11} &
Уравнение с параметром &
Решение не выполнено; требуется факторизация и подсчёт числа различных действительных корней при разных значениях параметра. &
не указан &
$a=0$ или $|a|>1$, кроме $a=-2$ \\

\hyperref[task:12]{12} &
Монеты и представление сумм &
Пункты а) и б) решены, но пункт в) оставлен без ответа, поэтому задание выполнено не полностью. &
а) да; б) нет; в) не указан &
а) да; б) нет; в) $3$ монеты \\

\hyperref[task:18]{18} &
Геометрия: касания и равнобедренный треугольник &
Намечены верные свойства центра окружности, но доказательство пункта а) не завершено и пункт б) не решён. &
не указан &
а) $\angle OAC=45^\circ$; б) $BC=5$ \\

\hyperref[task:20]{20} &
Правильная шестиугольная пирамида &
Выполнен только рисунок; доказательство параллельности и вычисление объёма отсутствуют. &
не указан &
а) $BC\parallel MK$; б) $72\sqrt3$ \\

\end{longtable}
\normalsize

\newpage
\thispagestyle{empty}
\vspace*{\fill}
\begin{center}
{\Huge\bfseries Разбор ошибок}
\end{center}
\vspace*{\fill}

\newpage
\phantomsection\label{task:3}
\section*{Задание 3}

\subsection*{Текст задания}
На рисунке изображены график функции $y=f(x)$ и касательная к нему в точке с абсциссой $x_0$. Найдите значение производной функции $f(x)$ в точке $x_0$.

\subsection*{Ваше решение}
В работе записано: «не знаю». Вычисления не выполнены.

\subsection*{Ваш ответ}
Не указан.

\subsection*{Ошибка}
\textcolor{red}{Ошибка:} решение не начато, поэтому не использован геометрический смысл производной.

\subsection*{Где именно возникает ошибка}
Последний корректный шаг отсутствует: после чтения условия вычислительный ход не начат.

Первый неполный шаг --- не определён угловой коэффициент изображённой касательной. Именно этот коэффициент и равен $f'(x_0)$.

\subsection*{Почему это неверно}
Производная функции в точке равна угловому коэффициенту касательной к графику функции в этой точке. Поэтому здесь не требуется восстанавливать формулу самой функции $f(x)$ и не требуется знать значение $x_0$.

На касательной отмечены две удобные точки:
\[
(-1;3)\quad\text{и}\quad(4;2).
\]
Если при переходе от первой точки ко второй абсцисса увеличивается на
\[
4-(-1)=5,
\]
то ордината уменьшается на
\[
2-3=-1.
\]
Следовательно, угловой коэффициент касательной равен отношению изменения ординаты к изменению абсциссы.

\subsection*{Ключевая идея}
\textcolor{blue!60!black}{Ключевая идея:} значение $f'(x_0)$ нужно снять непосредственно с касательной как её наклон.

\subsection*{Исправление}
Используем две отмеченные точки касательной:
\[
f'(x_0)=\frac{2-3}{4-(-1)}
       =\frac{-1}{5}
       =-\frac15.
\]

\subsection*{Полное правильное решение}
Касательная проходит через точки $(-1;3)$ и $(4;2)$. Её угловой коэффициент:
\[
k=\frac{\Delta y}{\Delta x}
  =\frac{2-3}{4-(-1)}
  =-\frac15.
\]
По геометрическому смыслу производной
\[
f'(x_0)=k.
\]
Значит,
\[
\boxed{f'(x_0)=-\frac15}.
\]

\subsection*{Проверка}
При увеличении $x$ на $5$ единиц значение $y$ на касательной уменьшается на $1$ единицу. Поэтому наклон действительно отрицательный и равен $-\frac15$, что согласуется с рисунком.

\subsection*{Итог}
\textcolor{teal!60!black}{Итог:} производная в точке касания равна
\[
-\frac15.
\]

\subsection*{Задача на закрепление}
К графику функции $y=g(x)$ в точке с абсциссой $x_1$ проведена касательная, проходящая через точки $(-2;5)$ и $(4;2)$. Найдите $g'(x_1)$.

\newpage
\phantomsection\label{task:7}
\section*{Задание 7}

\subsection*{Текст задания}
В июле 2030 года планируется взять кредит $182\,000$ рублей в банке на три года. Каждый январь долг увеличивается на $20\%$ по сравнению с концом предыдущего года. С февраля по июнь каждого года необходимо одним платежом выплатить часть долга. Найдите общую сумму платежей после полного погашения кредита, если кредит погашается тремя равными платежами.

\subsection*{Ваше решение}
В работе записано: «не знаю». Расчёт долга по годам не выполнен.

\subsection*{Ваш ответ}
Не указан.

\subsection*{Ошибка}
\textcolor{red}{Ошибка:} не составлена схема изменения долга после начисления процентов и очередного равного платежа.

\subsection*{Где именно возникает ошибка}
Последний корректный шаг отсутствует: после условия не введён размер равного платежа и не записан долг после первого года.

Первый неполный шаг --- не учтено, что перед каждым из трёх платежей остаток долга сначала увеличивается в $1{,}2$ раза.

\subsection*{Почему это неверно}
Проценты начисляются не на первоначальную сумму каждый год, а на фактический остаток долга на конец предыдущего года. Поэтому нельзя просто трижды прибавить $20\%$ к исходной сумме и затем разделить результат на три.

Нужно последовательно выполнять две операции:
\[
\text{остаток долга}\ \longrightarrow\ 1{,}2\cdot\text{остаток долга}
\longrightarrow\ \text{вычесть платёж}.
\]
Платежи равны, поэтому обозначим каждый из них через $x$.

\subsection*{Ключевая идея}
\textcolor{blue!60!black}{Ключевая идея:} записать остаток долга после каждого из трёх одинаковых платежей и потребовать, чтобы после третьего платежа долг стал равен нулю.

\subsection*{Исправление}
После начисления процентов в январе 2031 года долг равен
\[
182\,000\cdot1{,}2=218\,400.
\]
После первого платежа:
\[
218\,400-x.
\]
После начисления процентов в январе 2032 года и второго платежа:
\[
1{,}2(218\,400-x)-x.
\]
После начисления процентов в январе 2033 года третий платёж должен полностью закрыть долг:
\[
1{,}2\bigl(1{,}2(218\,400-x)-x\bigr)-x=0.
\]

\subsection*{Полное правильное решение}
Раскроем скобки:
\[
1{,}2\bigl(262\,080-1{,}2x-x\bigr)-x=0,
\]
\[
314\,496-2{,}64x-x=0,
\]
\[
314\,496=3{,}64x.
\]
Отсюда
\[
x=86\,400.
\]
Всего выполняются три одинаковых платежа, поэтому общая сумма:
\[
3\cdot86\,400=259\,200.
\]
Следовательно,
\[
\boxed{259\,200\text{ руб.}}
\]

\subsection*{Проверка}
Проверим долг по годам:
\[
182\,000\cdot1{,}2=218\,400,\qquad
218\,400-86\,400=132\,000;
\]
\[
132\,000\cdot1{,}2=158\,400,\qquad
158\,400-86\,400=72\,000;
\]
\[
72\,000\cdot1{,}2=86\,400,\qquad
86\,400-86\,400=0.
\]
После третьего платежа долг действительно полностью погашен.

\subsection*{Итог}
\textcolor{teal!60!black}{Итог:} каждый платёж равен $86\,400$ рублей, а общая сумма трёх платежей равна
\[
259\,200\text{ рублей}.
\]

\subsection*{Задача на закрепление}
В июле планируется взять кредит $210\,000$ рублей на два года. Каждый январь долг увеличивается на $10\%$ по сравнению с концом предыдущего года, а с февраля по июнь каждого года вносится один платёж. Кредит погашается двумя равными платежами. Найдите общую сумму платежей.

\newpage
\phantomsection\label{task:10}
\section*{Задание 10}

\subsection*{Текст задания}
Генератор имеет ЭДС $\mathcal E=50$ В и внутреннее сопротивление $r=1$ Ом. Сила тока в цепи определяется формулой
\[
I=\frac{\mathcal E}{r+R},
\]
а тепловая мощность, выделяемая в нагрузке, равна
\[
P=I^2R.
\]
Найдите максимально возможную мощность $P$.

\subsection*{Ваше решение}
Получено:
\[
P=\left(\frac{50}{1+R}\right)^2R
 =\frac{2500R}{(1+R)^2}.
\]
Далее выражение преобразовано, но максимальное значение мощности не найдено.

\subsection*{Ваш ответ}
Не указан.

\subsection*{Ошибка}
\textcolor{red}{Ошибка:} решение остановлено после получения правильной формулы мощности.

\subsection*{Где именно возникает ошибка}
Последний корректный шаг:
\[
P=\frac{2500R}{(1+R)^2}.
\]
Первый неполный шаг --- не выполнено сравнение знаменателя с числителем, которое даёт верхнюю границу для $P$.

\subsection*{Почему это неверно}
Сопротивление нагрузки положительно: $R>0$. Поэтому можно использовать очевидное неравенство
\[
(R-1)^2\ge 0.
\]
После раскрытия скобок:
\[
R^2-2R+1\ge 0,
\]
откуда
\[
(R+1)^2\ge 4R.
\]
Именно это позволяет оценить дробь
\[
\frac{R}{(R+1)^2}
\]
сверху и одновременно определить, когда достигается равенство.

\subsection*{Ключевая идея}
\textcolor{blue!60!black}{Ключевая идея:} не нужно перебирать значения $R$; достаточно получить точную верхнюю оценку для мощности и проверить условие равенства.

\subsection*{Исправление}
Из
\[
(R+1)^2\ge4R
\]
следует
\[
\frac{R}{(R+1)^2}\le\frac14.
\]
Поэтому
\[
P=2500\cdot\frac{R}{(R+1)^2}
\le 2500\cdot\frac14=625.
\]
Равенство достигается при $R=1$.

\subsection*{Полное правильное решение}
Подставляем $\mathcal E=50$ и $r=1$:
\[
I=\frac{50}{1+R}.
\]
Тогда
\[
P=I^2R
 =\left(\frac{50}{1+R}\right)^2R
 =\frac{2500R}{(1+R)^2}.
\]
Так как
\[
(R-1)^2\ge0,
\]
имеем
\[
(R+1)^2\ge4R.
\]
Следовательно,
\[
P\le625.
\]
При $R=1$ равенство выполняется, поэтому максимальная мощность действительно равна
\[
\boxed{625\text{ Вт}}.
\]

\subsection*{Проверка}
При $R=1$ Ом:
\[
I=\frac{50}{1+1}=25\text{ А}.
\]
Тогда
\[
P=25^2\cdot1=625\text{ Вт}.
\]
Полученное значение достигается, следовательно, это не только верхняя оценка, но и настоящий максимум.

\subsection*{Итог}
\textcolor{teal!60!black}{Итог:} максимальная тепловая мощность нагрузки равна
\[
625\text{ Вт}.
\]

\subsection*{Задача на закрепление}
Генератор имеет ЭДС $40$ В и внутреннее сопротивление $2$ Ом. Для нагрузки сопротивлением $R$ сила тока равна $\dfrac{40}{2+R}$, а мощность нагрузки равна $I^2R$. Найдите максимально возможную мощность нагрузки.

\newpage
\phantomsection\label{task:11}
\section*{Задание 11}

\subsection*{Текст задания}
Найдите все значения параметра $a$, при каждом из которых уравнение
\[
ax^4-2x^3+(a^3-a)x^2-(2a^2-2)x
=
-ax^3+2x^2-(a^3-a)x+2a^2-2
\]
имеет ровно два решения.

\subsection*{Ваше решение}
В работе записано: «не знаю». Преобразование уравнения не выполнено.

\subsection*{Ваш ответ}
Не указан.

\subsection*{Ошибка}
\textcolor{red}{Ошибка:} не выполнена факторизация уравнения и не проведён подсчёт различных действительных корней в зависимости от параметра.

\subsection*{Где именно возникает ошибка}
Последний корректный шаг отсутствует: решение не начато.

Первый неполный шаг --- не перенесены все члены в одну часть и не замечено разложение многочлена на три множителя.

\subsection*{Почему это неверно}
В задачах на число решений с параметром недостаточно просто решить уравнение при нескольких значениях $a$. Нужно получить структуру корней для всех возможных значений параметра и отдельно учесть:
\begin{itemize}
  \item исчезновение линейного корня при $a=0$;
  \item появление или исчезновение корней квадратного множителя;
  \item совпадение корней разных множителей;
  \item число \emph{различных} действительных решений.
\end{itemize}

После переноса правой части влево выражение раскладывается:
\[
(x+1)(ax-2)(x^2+a^2-1)=0.
\]

\subsection*{Ключевая идея}
\textcolor{blue!60!black}{Ключевая идея:} после факторизации нужно не просто выписать корни каждого множителя, а посчитать, сколько различных действительных чисел получается в каждом диапазоне значений $a$.

\subsection*{Исправление}
Рассматриваем множители:
\[
x+1=0,\qquad ax-2=0,\qquad x^2+a^2-1=0.
\]
Первый всегда даёт
\[
x=-1.
\]
Второй при $a\ne0$ даёт
\[
x=\frac2a.
\]
Третий имеет действительные корни тогда и только тогда, когда
\[
1-a^2\ge0,
\]
то есть $|a|\le1$.

\subsection*{Полное правильное решение}
После переноса всех членов в левую часть получаем
\[
(x+1)(ax-2)(x^2+a^2-1)=0.
\]

\textbf{Случай 1: $a=0$.}

Тогда
\[
(x+1)(-2)(x^2-1)=0.
\]
Так как
\[
x^2-1=(x-1)(x+1),
\]
различные корни:
\[
x=-1,\qquad x=1.
\]
Их ровно два, поэтому
\[
a=0
\]
подходит.

\textbf{Случай 2: $0<|a|<1$.}

Квадратный множитель даёт два различных корня
\[
x=\pm\sqrt{1-a^2}.
\]
Они лежат строго между $-1$ и $1$. Кроме них есть корень $x=-1$.

Также существует корень
\[
x=\frac2a.
\]
Поскольку $0<|a|<1$, имеем
\[
\left|\frac2a\right|>2,
\]
поэтому этот корень не совпадает ни с $-1$, ни с корнями квадратного множителя.

Получается четыре различных действительных корня. Такой параметр не подходит.

\textbf{Случай 3: $|a|=1$.}

Если $a=1$, корни:
\[
-1,\quad 2,\quad 0.
\]
Если $a=-1$, корни:
\[
-1,\quad -2,\quad 0.
\]
В обоих случаях различных корней три. Значения $a=\pm1$ не подходят.

\textbf{Случай 4: $|a|>1$.}

Квадратный множитель
\[
x^2+a^2-1
\]
действительных корней не имеет. Остаются
\[
x=-1,\qquad x=\frac2a.
\]
Обычно это два различных корня. Они совпадают только при
\[
\frac2a=-1,
\]
откуда
\[
a=-2.
\]
При $a=-2$ остаётся только один различный действительный корень, поэтому это значение нужно исключить.

Итак,
\[
\boxed{a=0\quad\text{или}\quad |a|>1,\ a\ne-2}.
\]
В интервальной записи:
\[
\boxed{a\in(-\infty,-2)\cup(-2,-1)\cup\{0\}\cup(1,\infty)}.
\]

\subsection*{Проверка}
Проверим характерные значения:
\[
a=0 \Rightarrow \{-1,1\},
\]
то есть два корня.

При $a=2$ квадратный множитель корней не даёт, а линейные множители дают $-1$ и $1$ --- два корня.

При $a=-2$ оба линейных множителя дают один и тот же корень $-1$, поэтому остаётся только одно различное решение. Исключение $a=-2$ необходимо.

\subsection*{Итог}
\textcolor{teal!60!black}{Итог:} ровно два различных действительных решения получаются при
\[
a\in(-\infty,-2)\cup(-2,-1)\cup\{0\}\cup(1,\infty).
\]

\subsection*{Задача на закрепление}
Найдите все значения параметра $b$, при которых уравнение
\[
(x-2)(bx+1)(x^2+b^2-4)=0
\]
имеет ровно два различных действительных корня.

\newpage
\phantomsection\label{task:12}
\section*{Задание 12}

\subsection*{Текст задания}
Есть $16$ монет по $2$ рубля и $29$ монет по $5$ рублей.

а) Можно ли выбрать несколько монет так, чтобы сумма была равна $175$ рублям?

б) Можно ли получить сумму $176$ рублей?

в) Какое наименьшее количество монет по $1$ рублю нужно добавить, чтобы можно было получить любую целую сумму от $1$ до $180$ рублей включительно?

\subsection*{Ваше решение}
Для пункта а) верно замечено, что общая стоимость всех имеющихся монет равна
\[
16\cdot2+29\cdot5=177,
\]
и из неё можно убрать одну монету достоинством $2$ рубля, получив $175$ рублей.

Для пункта б) сделан вывод, что получить $176$ рублей невозможно.

Пункт в) оставлен без решения.

\subsection*{Ваш ответ}
а) да; б) нет; в) не указан.

\subsection*{Ошибка}
\textcolor{red}{Ошибка:} задание выполнено не полностью: отсутствует решение пункта в).

\subsection*{Где именно возникает ошибка}
Последний корректный шаг --- вывод в пункте б), что сумма $176$ рублей недостижима.

Первый неполный шаг --- после пункта б) не определено минимальное число однорублёвых монет и не доказано, что найденного количества достаточно для получения \emph{каждой} суммы от $1$ до $180$.

\subsection*{Почему это неверно}
В пункте в) требуется выполнить две отдельные части доказательства:
\begin{enumerate}
  \item показать, что меньше некоторого количества однорублёвых монет быть не может;
  \item показать, что найденного количества действительно достаточно для всех сумм из указанного диапазона.
\end{enumerate}
Одной проверки общей стоимости монет недостаточно: кроме суммы $180$ нужно обеспечить отсутствие «пропусков» среди меньших сумм.

\subsection*{Ключевая идея}
\textcolor{blue!60!black}{Ключевая идея:} сначала получить нижнюю границу по общей стоимости, затем доказать непрерывность набора достижимых сумм.

\subsection*{Исправление}
Стоимость исходного набора:
\[
177\text{ рублей}.
\]
Чтобы вообще можно было получить $180$ рублей, необходимо добавить не менее
\[
180-177=3
\]
однорублёвых монет.

Остаётся доказать, что трёх монет достаточно.

\subsection*{Полное правильное решение}
\textbf{Пункт а).}

Общая стоимость всех монет:
\[
16\cdot2+29\cdot5=32+145=177.
\]
Если не брать одну двухрублёвую монету, получим
\[
177-2=175.
\]
Следовательно, ответ: да.

\textbf{Пункт б).}

Пусть выбрано $x$ двухрублёвых монет и $y$ пятирублёвых. Тогда
\[
2x+5y=176,
\]
где
\[
0\le x\le16,\qquad 0\le y\le29.
\]
Так как $176$ и $2x$ чётные, число $5y$ должно быть чётным, значит $y$ чётно. Поэтому
\[
y\le28.
\]
Тогда максимально возможная сумма при чётном $y$ равна
\[
2\cdot16+5\cdot28=32+140=172<176.
\]
Следовательно, получить $176$ рублей невозможно.

\textbf{Пункт в).}

Исходные монеты дают всего $177$ рублей. Поэтому для суммы $180$ рублей нужно добавить как минимум три однорублёвые монеты:
\[
m\ge3.
\]

Покажем, что $m=3$ достаточно.

Три однорублёвые монеты позволяют получить любую сумму от $0$ до $3$ рублей.

Добавляем шестнадцать двухрублёвых монет. Поскольку очередная монета имеет достоинство $2$ рубля, а до её добавления уже можно получить все суммы подряд начиная с нуля, пропусков не возникает. После использования всех шестнадцати двухрублёвых монет можно получить любую сумму от
\[
0
\]
до
\[
3+16\cdot2=35.
\]

Теперь добавляем пятирублёвые монеты. Первая такая монета равна $5$, а до этого доступны все суммы от $0$ до $35$. Поэтому с ней получаются все суммы от $5$ до $40$, и вместе с прежними суммами --- весь диапазон от $0$ до $40$ без пропусков.

Каждая следующая пятирублёвая монета снова расширяет непрерывный диапазон ещё на $5$ рублей. После всех $29$ пятирублёвых монет верхняя граница станет
\[
3+16\cdot2+29\cdot5=180.
\]
Значит, с тремя однорублёвыми монетами можно получить любую целую сумму от $1$ до $180$.

Следовательно,
\[
\boxed{3}
\]
--- минимальное количество.

\subsection*{Проверка}
Меньше трёх монет не хватит уже по общей стоимости:
\[
177+2=179<180.
\]
Три монеты дают общую стоимость ровно
\[
177+3=180,
\]
а приведённое выше построение показывает, что доступны не только $180$ рублей, но и все промежуточные целые суммы.

\subsection*{Итог}
\textcolor{teal!60!black}{Итог:} ответы:
\[
\text{а) да;\qquad б) нет;\qquad в) }3\text{ монеты по 1 рублю}.
\]

\subsection*{Задача на закрепление}
Есть $12$ монет по $2$ рубля и $20$ монет по $5$ рублей. Какое наименьшее количество монет по $1$ рублю нужно добавить, чтобы можно было получить любую целую сумму от $1$ до $130$ рублей включительно?

\newpage
\phantomsection\label{task:18}
\section*{Задание 18}

\subsection*{Текст задания}
Окружность с центром $O$ касается боковых сторон $AB$ и $BC$ равнобедренного остроугольного треугольника $ABC$ и его высоты $CH$.

а) Докажите, что $\angle OAC=45^\circ$.

б) Найдите $BC$, если $BO=1$ и $AC=6$.

\subsection*{Ваше решение}
Верно отмечено, что из касания окружности сторон $AB$ и $BC$ следует расположение центра $O$ на биссектрисе угла $ABC$. Также намечено использование свойства биссектрисы, проведённой из вершины равнобедренного треугольника, и равенства $OA=OC$.

Дальнейшее доказательство не завершено, пункт б) не решён.

\subsection*{Ваш ответ}
Не указан.

\subsection*{Ошибка}
\textcolor{red}{Ошибка:} решение остановлено после выбора правильных геометрических свойств; необходимые углы не связаны между собой, а числовая часть задания не выполнена.

\subsection*{Где именно возникает ошибка}
Последний корректный шаг --- центр окружности $O$ лежит на биссектрисе угла $ABC$. Поскольку $AB=BC$, эта биссектриса одновременно является медианой и высотой к основанию $AC$, то есть лежит на серединном перпендикуляре к $AC$. Поэтому действительно
\[
OA=OC.
\]

Первый неполный шаг --- после этого не найден угол $\angle ACO$ и не использовано равенство $OA=OC$ для получения требуемого угла $\angle OAC$.

\subsection*{Почему это неверно}
Само равенство $OA=OC$ ещё не завершает доказательство: оно только показывает, что треугольник $AOC$ равнобедренный. Чтобы определить его углы, нужно использовать третье касание --- к высоте $CH$.

Поскольку окружность касается прямых $BC$ и $CH$, её центр равноудалён от этих прямых. Значит, $CO$ является биссектрисой угла $BCH$.

Именно сочетание двух биссектрис --- при вершинах $B$ и $C$ --- позволяет получить угол $45^\circ$.

\subsection*{Ключевая идея}
\textcolor{blue!60!black}{Ключевая идея:} центр окружности, касающейся двух пересекающихся прямых, лежит на биссектрисе угла между ними. Здесь это свойство нужно применить дважды.

\subsection*{Исправление}
Обозначим
\[
\angle BAC=\angle ACB=\alpha,
\]
так как $AB=BC$.

Поскольку $CH\perp AB$,
\[
\angle ACH=90^\circ-\alpha.
\]

Угол при вершине $B$ равен
\[
\angle ABC=180^\circ-2\alpha.
\]
В прямоугольном треугольнике $BCH$:
\[
\angle BCH=90^\circ-\angle ABC
=90^\circ-(180^\circ-2\alpha)
=2\alpha-90^\circ.
\]
Так как $CO$ --- биссектриса угла $BCH$,
\[
\angle HCO=\alpha-45^\circ.
\]
Следовательно,
\[
\angle ACO
=
\angle ACH+\angle HCO
=
(90^\circ-\alpha)+(\alpha-45^\circ)
=
45^\circ.
\]
А так как $OA=OC$, то
\[
\angle OAC=\angle ACO=45^\circ.
\]

\subsection*{Полное правильное решение}
\textbf{Пункт а).}

Окружность касается сторон $AB$ и $BC$, поэтому центр $O$ лежит на биссектрисе угла $ABC$.

Треугольник $ABC$ равнобедренный, $AB=BC$. Биссектриса из вершины $B$ является также медианой и высотой к основанию $AC$. Значит, эта прямая --- серединный перпендикуляр к $AC$. Поскольку $O$ лежит на ней,
\[
OA=OC.
\]

Пусть
\[
\angle BAC=\angle ACB=\alpha.
\]
Так как $CH\perp AB$,
\[
\angle ACH=90^\circ-\alpha.
\]
Далее,
\[
\angle ABC=180^\circ-2\alpha.
\]
В прямоугольном треугольнике $BCH$:
\[
\angle BCH=2\alpha-90^\circ.
\]
Окружность касается $BC$ и $CH$, поэтому $CO$ --- биссектриса угла $BCH$:
\[
\angle HCO=\alpha-45^\circ.
\]
Тогда
\[
\angle ACO
=
(90^\circ-\alpha)+(\alpha-45^\circ)
=
45^\circ.
\]
Из $OA=OC$ следует равенство углов при основании треугольника $AOC$:
\[
\boxed{\angle OAC=45^\circ}.
\]

\textbf{Пункт б).}

Снова обозначим
\[
\angle BAC=\alpha.
\]
Из пункта а)
\[
\angle OAC=45^\circ.
\]
Поэтому
\[
\angle BAO=\alpha-45^\circ.
\]

Так как $BO$ --- биссектриса угла $ABC$,
\[
\angle ABO
=
\frac{180^\circ-2\alpha}{2}
=
90^\circ-\alpha.
\]
Следовательно, в треугольнике $AOB$
\[
\angle AOB
=
180^\circ-(\alpha-45^\circ)-(90^\circ-\alpha)
=
135^\circ.
\]

По теореме синусов:
\[
\frac{AB}{\sin135^\circ}
=
\frac{BO}{\sin(\alpha-45^\circ)}.
\]
Так как $BO=1$,
\[
AB=
\frac{\sin135^\circ}{\sin(\alpha-45^\circ)}
=
\frac{1}{\sin\alpha-\cos\alpha}.
\]

Теперь используем $AC=6$. Высота из вершины $B$ в равнобедренном треугольнике делит основание пополам, поэтому в образовавшемся прямоугольном треугольнике
\[
\cos\alpha=\frac{3}{AB},
\]
то есть
\[
AB=\frac{3}{\cos\alpha}.
\]
Приравниваем два выражения для $AB$:
\[
\frac{3}{\cos\alpha}
=
\frac{1}{\sin\alpha-\cos\alpha}.
\]
Отсюда
\[
3\sin\alpha-3\cos\alpha=\cos\alpha,
\]
\[
3\sin\alpha=4\cos\alpha,
\]
\[
\tan\alpha=\frac43.
\]
Так как $\alpha$ острый, получаем стандартное отношение сторон прямоугольного треугольника $3$--$4$--$5$:
\[
\cos\alpha=\frac35.
\]
Тогда
\[
AB=\frac{3}{3/5}=5.
\]
Поскольку $AB=BC$,
\[
\boxed{BC=5}.
\]

\subsection*{Проверка}
При $BC=AB=5$ и $AC=6$ половина основания равна $3$. Высота равнобедренного треугольника равна
\[
\sqrt{5^2-3^2}=4.
\]
Получается треугольник со сторонами $5,5,6$, что согласуется с
\[
\tan\alpha=\frac43.
\]
Все углы острые, как и требуется в условии.

\subsection*{Итог}
\textcolor{teal!60!black}{Итог:} доказано, что
\[
\angle OAC=45^\circ,
\]
а при $BO=1$ и $AC=6$
\[
BC=5.
\]

\subsection*{Задача на закрепление}
Окружность с центром $O$ касается боковых сторон $AB$ и $BC$ равнобедренного остроугольного треугольника $ABC$ и его высоты $CH$. Известно, что $BO=2$ и $AC=12$. Найдите длину $BC$.

\newpage
\phantomsection\label{task:20}
\section*{Задание 20}

\subsection*{Текст задания}
В правильной шестиугольной пирамиде $SABCDEF$ точки $M$ и $K$ --- середины боковых рёбер $SA$ и $SD$ соответственно.

а) Докажите, что прямые $BC$ и $MK$ параллельны.

б) Найдите объём пирамиды $MABF$, если $AB=12$, а угол между плоскостью $BMC$ и плоскостью основания пирамиды $SABCDEF$ равен $30^\circ$.

\subsection*{Ваше решение}
Выполнен рисунок и отмечены точки $M$ и $K$ как середины боковых рёбер. Доказательство пункта а) и вычисления пункта б) отсутствуют.

\subsection*{Ваш ответ}
Не указан.

\subsection*{Ошибка}
\textcolor{red}{Ошибка:} решение не доведено дальше построения рисунка.

\subsection*{Где именно возникает ошибка}
Последний корректный шаг --- корректно отмечены середины рёбер $SA$ и $SD$.

Первый неполный шаг --- не применено свойство средней линии треугольника $SAD$, из которого сразу получается направление $MK$. В пункте б) не построено линейное представление двугранного угла между плоскостями.

\subsection*{Почему это неверно}
Для пункта а) нужно связать пространственный отрезок $MK$ со стороной основания: точки $M$ и $K$ лежат в одной грани $SAD$, поэтому работает обычная теорема о средней линии треугольника.

Для пункта б) угол между плоскостями нельзя подменять произвольным углом между линиями этих плоскостей. Плоскости $BMC$ и основания пересекаются по прямой $BC$. Поэтому нужно взять по одной прямой в каждой плоскости, перпендикулярной $BC$ в одной и той же точке.

\subsection*{Ключевая идея}
\textcolor{blue!60!black}{Ключевая идея:} в пункте а) использовать среднюю линию треугольника $SAD$, а в пункте б) рассмотреть сечение, перпендикулярное общей прямой $BC$ двух плоскостей.

\subsection*{Исправление}
В треугольнике $SAD$ точки $M$ и $K$ --- середины $SA$ и $SD$. Поэтому
\[
MK\parallel AD.
\]
В правильном шестиугольнике $ABCDEF$
\[
AD\parallel BC.
\]
Следовательно,
\[
MK\parallel BC.
\]

Для пункта б) обозначим через $O$ центр основания, а высоту пирамиды через
\[
SO=h.
\]
Так как $M$ --- середина $SA$, расстояние от $M$ до плоскости основания равно
\[
\frac h2.
\]

\subsection*{Полное правильное решение}
\textbf{Пункт а).}

В треугольнике $SAD$ отрезок $MK$ соединяет середины двух сторон:
\[
M\in SA,\qquad K\in SD.
\]
По теореме о средней линии треугольника
\[
MK\parallel AD.
\]
В правильном шестиугольнике диагональ $AD$ параллельна стороне $BC$. Поэтому
\[
\boxed{MK\parallel BC}.
\]

\textbf{Пункт б).}

Пусть $O$ --- центр правильного шестиугольника $ABCDEF$, а
\[
SO=h.
\]
Расположим основание мысленно так, чтобы $AD$ и $BC$ были горизонтальны. В правильном шестиугольнике со стороной $a=AB=12$ расстояние от центра $O$ до стороны $BC$ равно
\[
\frac{\sqrt3}{2}a.
\]

Проекция точки $M$, являющейся серединой $SA$, на основание есть середина отрезка $OA$. Обозначим эту проекцию через $P$. Для правильного шестиугольника прямая $PB$ перпендикулярна $BC$, причём
\[
PB=\frac{\sqrt3}{2}a.
\]
Так как вертикальная составляющая отрезка $MB$ равна расстоянию от $M$ до основания,
\[
MP=\frac h2.
\]
Кроме того, $MB\perp BC$. Поэтому угол между $MB$ и его проекцией $BP$ на основание как раз равен углу между плоскостью $BMC$ и плоскостью основания:
\[
\angle MBP=30^\circ.
\]

В прямоугольном треугольнике $MBP$:
\[
\tan30^\circ=\frac{MP}{PB}
=
\frac{h/2}{(\sqrt3/2)a}
=
\frac{h}{\sqrt3\,a}.
\]
Так как
\[
\tan30^\circ=\frac1{\sqrt3},
\]
получаем
\[
\frac1{\sqrt3}=\frac{h}{\sqrt3\,a},
\]
откуда
\[
h=a=12.
\]

Теперь найдём площадь основания пирамиды $MABF$. Её основание --- треугольник $ABF$ в плоскости правильного шестиугольника.

В правильном шестиугольнике
\[
BF=\sqrt3\,a,
\]
а расстояние от точки $A$ до прямой $BF$ равно
\[
\frac a2.
\]
Поэтому
\[
S_{ABF}
=
\frac12\cdot \sqrt3\,a\cdot\frac a2
=
\frac{\sqrt3\,a^2}{4}.
\]
При $a=12$:
\[
S_{ABF}
=
\frac{\sqrt3\cdot144}{4}
=
36\sqrt3.
\]

Высота пирамиды $MABF$ равна расстоянию от $M$ до плоскости основания:
\[
\frac h2=6.
\]
Следовательно,
\[
V_{MABF}
=
\frac13 S_{ABF}\cdot6
=
\frac13\cdot36\sqrt3\cdot6
=
72\sqrt3.
\]
Итак,
\[
\boxed{V_{MABF}=72\sqrt3}.
\]

\subsection*{Проверка}
При $h=12$ расстояние от $M$ до основания равно $6$. Для $a=12$
\[
PB=6\sqrt3.
\]
Поэтому
\[
\tan\angle MBP=\frac{6}{6\sqrt3}=\frac1{\sqrt3},
\]
значит,
\[
\angle MBP=30^\circ,
\]
что совпадает с условием.

\subsection*{Итог}
\textcolor{teal!60!black}{Итог:}
\[
BC\parallel MK,
\qquad
V_{MABF}=72\sqrt3.
\]

\subsection*{Задача на закрепление}
В правильной шестиугольной пирамиде $SABCDEF$ точка $M$ --- середина ребра $SA$. Сторона основания равна $6$, а угол между плоскостью $BMC$ и плоскостью основания равен $30^\circ$. Найдите объём пирамиды $MABF$.

\vfill
\begin{center}
\small
Лёвин Артём Александрович эксклюзивно для Матвея Горбачева
\end{center}

\end{document}
